\begin{exsol}{11}{\MR{} Food for Thought 1.3}{18}{Exam}{A1}
\asked Given the 1991 Treasury-auction scandal at Salomon Brothers, judge whether the episode shows
that the firm's incentive scheme, as described in Chapter~1, was fundamentally flawed.

\method The Salomon case (\MR{} Ch.~1, pp.~9--12): bonus close to an individual piece rate on
attributed profit (pp.~10--11), its known side effects (p.~11), and the May 1990 reform paying part
of the bonus in stock held in trust for five years (pp.~11--12). The patterns of pp.~16--17:
authority where the information is, incentives and delegation as complements. Complementarity and
coherent systems (Ch.~4, pp.~108--113). The question requires separating three levers:
\emph{incentives} (what is rewarded), \emph{selection} (who is attracted and promoted) and
\emph{control} (what actions are allowed, monitored and punished).

\solution
\subparagraph{Incentives: a real defect, but not the one the reform targeted.} The bonus rewarded
measured profit attributed to a desk (p.~10). Two costs were missing from that measure: legal
and regulatory risk, and damage to the firm's reputation, a firm-wide asset. A trader who corners
an auction collects the measured profit; the firm bears the legal penalty, the lost customers and
the fall in share price. When the downside is mostly the firm's and the upside largely the
trader's, high-powered pay on measured profit rewards excessive risk, including the risk of
breaking rules. In that precise sense the scheme was flawed: it paid for a measure that omitted
important costs. \MR{} had already listed related side effects --- little cooperation, focus on
own accounts, ``stealing'' other departments' profits (p.~11).

The 1990 reform aimed at this problem: stock held in trust for five years ties each bonus to the
firm's long-run value (p.~11), and the scandal did cut that value sharply. But stock pay is a
\emph{group} incentive. One managing director owns a tiny fraction of the firm, so his share of
the firm-wide loss is negligible next to the personal gain from a large trading profit. Stock
ownership can encourage cooperation and a long-run view across the whole staff; it cannot deter
one individual's rule-breaking. The reform was also barely in place when the violations began
(February 1991). The episode therefore does not refute the reform's logic, but it shows its limits.

\subparagraph{Selection and culture.} Profits were ``the sole source of power, prestige, status,
and respect'' (p.~10), and the pay system attracted individualistic, competitive risk-takers
(p.~10). One of the theories offered for the violation --- a competitive need to beat the
regulators rather than greed --- is a selection explanation: the marginal incentive may matter
less than \emph{who} the system attracts and promotes. \MR{} expected the reform to change the type
of people Salomon attracted only ``in the long run'' (p.~11).

\subparagraph{Control.} Senior executives learned of the violations and did not act. This is a
failure of monitoring and governance, not of the trader's pay formula. Salomon gave traders
enormous discretion because they held the information (p.~17), and the book stresses that such
delegation needs incentives to go with it. The case adds the other half: discretion over actions
with catastrophic, firm-wide downside also needs limits, compliance and reporting that someone is
paid --- and willing --- to enforce. The executives' own incentives pointed the other way:
reporting promptly would have meant immediate costs for them.

\subparagraph{Verdict.} Not fundamentally flawed as a device to elicit effort and use traders'
local information --- that was the scheme's purpose and it achieved it (p.~11). Flawed as a
\emph{system}: profit-based pay, a culture that ranked people by profit, selection of aggressive
risk-takers and weak controls were mutually reinforcing, a coherent pattern in the sense of
pp.~108--113. Because the parts are complements, changing one of them (adding stock to the bonus)
could not move the whole system; the fix had to reach controls, leadership and culture together.

\interp Judging an incentive scheme by one scandal confuses three levers. The right question is
which lever failed. Here all three contributed, and the case is a warning against reading pay
alone as ``the incentive system'': incentives, selection and control form one design.

\verify Counterfactual tests. With the same bonus but strong controls (position limits, compliance
checks, senior managers who report violations), the violation would have been far less likely: so
pay alone is not sufficient to explain it. With flat salaries, the 1980s trading profits and the
use of traders' information would probably have been lost (pp.~10--11): so removing incentive pay
is not the fix. The market evidence points the same way: the reform raised the share price by
$7$--$10\%$ (p.~12), so investors valued its incentive logic.

\answer{No, not fundamentally: the scheme did what it was designed to do, eliciting effort and
using traders' information. But it was flawed in what it measured (attributed profit that ignored
legal and reputational risk), and the 1990 stock reform is a group incentive that cannot deter one
individual. The scandal mainly exposes failures of selection, culture and control, which formed a
coherent system with the pay scheme; fixing it required changing all of them together.}
\end{exsol}
